% Vietnamese translation for OI Public Library
% Source-PDF: 比赛 CONTEST/2021“MINIEYE杯”中国大学生算法设计超级联赛/2021“MINIEYE杯”中国大学生算法设计超级联赛（4）/2021“MINIEYE杯”中国大学生算法设计超级联赛（4）-题目集.pdf
\documentclass[11pt]{article}
\usepackage{oipl-vn}

\title{Đề bài trận thứ tư HDU Multi-University 2021}
\author{}
\date{}

\begin{document}
\maketitle

\origtitle{Super League of Chinese College Students Algorithm Design 2021, Contest 4}
\sourcepath{contest/hdu-2021-minieye-04-problems.pdf}

\section{Problem A. Calculus}

\paragraph{Tệp vào:} standard input
\paragraph{Tệp ra:} standard output
\paragraph{Giới hạn bộ nhớ:} 256 megabytes

Mùa hè này, ZXyang đã quá mệt mỏi khi làm các bài của những cuộc thi
Multi-University. Vì vậy cậu quyết định tham dự Kỳ thi tuyển sinh sau đại học
thống nhất toàn quốc. Hôm nay, cậu thấy một bài toán về chuỗi.

Gọi \(S(x)\) là một hàm có biến độc lập là \(x\). \(S(x)\) có thể được biểu
diễn bởi công thức sau:
\[
  f(x)=\sum_{i=1}^{n} f_i(x),
  \qquad
  S(x)=\sum_{j=1}^{x} f(j).
\]

\(f_i(x)\) là một hàm có biến độc lập là \(x\). Hơn nữa, \(f_i(x)\) thuộc tập
hàm \(F\):
\[
  F=\left\{
  C,\frac{C}{x},C\sin x,C\cos x,
  \frac{C}{\sin x},\frac{C}{\cos x},Cx,C^x
  \right\}.
\]

\(C\) là một hằng số nguyên trong khoảng từ \(0\) đến \(10^9\).

ZXyang muốn biết \(S(x)\) có hội tụ hay không. \(S(x)\) hội tụ khi và chỉ khi
\[
  \lim_{x\to\infty} S(x)=c,
\]
trong đó \(c\) là một hằng số.

\subsection*{Dữ liệu vào}

Dòng đầu tiên chứa một số nguyên \(t\) (\(1\le t\le 10^4\)) là số bộ dữ liệu.

Dòng đầu tiên và cũng là dòng duy nhất của mỗi bộ dữ liệu chứa một chuỗi
\(s\) (\(1\le |s|\le 100\)), biểu thị công thức của \(f(x)\). Phân số được
viết dưới dạng \texttt{a/b}. \(C^x\) được viết dưới dạng \verb|C^x|.
Bảo đảm rằng hằng số \(C\) sẽ không bị lược bỏ khi \(C=1\). \(f(x)\) gồm các
hàm trong \(F\) nối với nhau bằng dấu \texttt{+}.

\subsection*{Dữ liệu ra}

Với mỗi bộ dữ liệu, in \texttt{YES} trên một dòng nếu \(S(x)\) là một dãy hội
tụ, hoặc in \texttt{NO} nếu không.

\subsection*{Ví dụ}

\paragraph{standard input}
\begin{verbatim}
2
1sinx+0cosx+3x+6/sinx
0
\end{verbatim}

\paragraph{standard output}
\begin{verbatim}
NO
YES
\end{verbatim}

\section{Problem B. Kanade Loves Maze Designing}

\paragraph{Tệp vào:} standard input
\paragraph{Tệp ra:} standard output
\paragraph{Giới hạn bộ nhớ:} 256 megabytes

Kanade đang thiết kế một trò chơi nhỏ. Đó là một trò chơi giải đố yêu cầu
người chơi thoát khỏi mê cung. Người chơi cũng nên thu thập càng nhiều loại
nguyên tố càng tốt để đạt điểm cao hơn.

Ở chế độ dễ, mê cung có thể được mô tả như một cây. Có \(n\) giao lộ và
\(n-1\) lối đi vô hướng làm cho \(n\) giao lộ liên thông. Các giao lộ được
đánh số bằng các số nguyên từ \(1\) đến \(n\). Trên mỗi giao lộ đặt đúng một
nguyên tố. Loại nguyên tố đặt tại giao lộ \(i\) được ký hiệu bởi một số
nguyên \(c_i\) trong khoảng \([1,n]\).

Để đánh giá độ khó của mê cung, Kanade muốn biết có bao nhiêu loại nguyên tố
xuất hiện trên \(p(u,v)\) với mọi hai số nguyên \(u,v\in[1,n]\). \(p(u,v)\)
chỉ đường đi đơn từ giao lộ \(u\) đến giao lộ \(v\) trong mê cung.

\subsection*{Dữ liệu vào}

Dòng đầu tiên chứa một số nguyên \(T\) (\(1\le T\le 10\)), biểu thị số bộ dữ
liệu.

Với mỗi bộ dữ liệu, dòng đầu tiên chứa một số nguyên \(n\)
(\(2\le n\le 2000\)), biểu thị số giao lộ trong mê cung.

Dòng thứ hai chứa \(n-1\) số nguyên \(p_2,p_3,\ldots,p_n\) (\(i<p_i\)), biểu
thị giao lộ thứ \(i\) được nối với giao lộ thứ \(p_i\) bằng một lối đi.

Dòng thứ ba chứa \(n\) số nguyên \(c_1,c_2,\ldots,c_n\) (\(1\le c_i\le n\)),
biểu thị loại nguyên tố đặt tại giao lộ \(i\) là \(c_i\).

Bảo đảm rằng với tất cả các bộ dữ liệu,
\[
  \sum n\le 5000.
\]

\subsection*{Dữ liệu ra}

Với mỗi bộ dữ liệu, in \(n\) dòng. Mỗi dòng chứa hai số nguyên. Gọi \(a_{i,j}\)
là số loại nguyên tố xuất hiện trên \(p(i,j)\). Đặt
\[
  f(i,x)=\sum_{j=1}^{n} a_{i,j}x^{j-1}.
\]

Ở dòng thứ \(i\), hãy in \(f(i,19560929)\bmod (10^9+7)\) và
\(f(i,19560929)\bmod (10^9+9)\), cách nhau bởi một dấu cách.

\subsection*{Ví dụ}

\paragraph{standard input}
\begin{verbatim}
1
6
1 1 2 2 3
1 1 4 5 1 4
\end{verbatim}

\paragraph{standard output}
\begin{verbatim}
495644981 518101442
495644981 518101442
397599492 896634980
612255048 326063507
495644981 518101442
397599492 896634980
\end{verbatim}

\subsection*{Ghi chú}

Gọi \(A=(a_{i,j})_{n\times n}\), trong ví dụ ta có
\[
  A=
  \begin{bmatrix}
  1&1&2&2&1&2\\
  1&1&2&2&1&2\\
  2&2&1&3&2&1\\
  2&2&3&1&2&3\\
  1&1&2&2&1&2\\
  2&2&1&3&2&1
  \end{bmatrix}.
\]

\section{Problem C. Cycle Binary}

\paragraph{Tệp vào:} standard input
\paragraph{Tệp ra:} standard output
\paragraph{Giới hạn bộ nhớ:} 512 megabytes

Ta có thể viết lại một chuỗi nhị phân, tức là chuỗi chỉ gồm các ký tự \(0\)
hoặc \(1\), \(s\) dưới dạng \(kp+p'\). Ở đây \(p'\) là một tiền tố của \(p\)
(\(p'\) có thể rỗng). \(+p'\) nghĩa là nối \(p'\) vào cuối \(kp\). \(kp\)
nghĩa là nối \(k\) bản sao của \(p\), trong đó \(p\) là đơn vị của \(s\).

Đơn vị của một chuỗi nhị phân \(s\) là chuỗi không rỗng \(p\) làm cho \(k\)
lớn nhất khi bỏ qua phần \(p'\).

Ta định nghĩa \(v(s)\) là giá trị của \(s\), tức là \(k\) trong mô tả trên.
Ví dụ, \(v(01001001)=2\), vì \(01001001\) có thể được viết thành
\(2(010)+(01)\), và \(v(11111)=5\) vì \(11111\) có thể được viết thành
\(5(1)\) và \(p'\) rỗng.

Cho \(n\), nhiệm vụ của bạn là tính tổng giá trị của tất cả các chuỗi nhị
phân có độ dài đúng bằng \(n\). Vì đáp án có thể rất lớn, chỉ cần in đáp án
theo modulo \(998244353\).

\subsection*{Dữ liệu vào}

Dòng đầu tiên chứa một số nguyên dương \(T\) (\(1\le T\le 100\)), biểu thị có
\(T\) bộ dữ liệu.

Mỗi bộ dữ liệu chứa một số nguyên dương \(n\) (\(1\le n\le 10^9\)) trên một
dòng.

Bảo đảm rằng
\[
  \sum n\le 10^{10}.
\]

\subsection*{Dữ liệu ra}

Với mỗi bộ dữ liệu, in một số nguyên trên một dòng, biểu thị đáp án theo
modulo \(998244353\).

\subsection*{Ví dụ}

\paragraph{standard input}
\begin{verbatim}
5
1
2
3
114
514
\end{verbatim}

\paragraph{standard output}
\begin{verbatim}
2
6
12
954037435
530871613
\end{verbatim}

\subsection*{Ghi chú}

Với \(n=3\),
\[
\begin{aligned}
 &v(000)+v(001)+v(010)+v(011)+v(100)+v(101)+v(110)+v(111)\\
 &=3+1+1+1+1+1+1+3\\
 &=12.
\end{aligned}
\]

\section{Problem D. Display Substring}

\paragraph{Tệp vào:} standard input
\paragraph{Tệp ra:} standard output
\paragraph{Giới hạn bộ nhớ:} 256 megabytes

Khi hiển thị một chuỗi trên màn hình LED, việc hiển thị các ký tự khác nhau
có thể tiêu thụ các mức năng lượng khác nhau. Cho một chuỗi \(S\) độ dài
\(n\) gồm các chữ cái tiếng Anh thường và mức tiêu thụ năng lượng của từng
chữ cái. Giả sử mức tiêu thụ năng lượng của một chuỗi là tổng mức tiêu thụ
năng lượng của tất cả các ký tự trong chuỗi đó. Hãy tìm mức tiêu thụ năng
lượng nhỏ thứ \(k\) trong tất cả các xâu con khác nhau của \(S\).

Chú ý rằng xâu con của một chuỗi \(S\) là một dãy con liên tiếp của \(S\). Ta
nói hai chuỗi \(S\) và \(T\) là khác nhau nếu độ dài của chúng khác nhau hoặc
tồn tại ít nhất một vị trí \(i\) sao cho \(S[i]\ne T[i]\).

\subsection*{Dữ liệu vào}

Dòng đầu tiên chứa một số nguyên \(T\) (\(1\le T\le 2\times 10^3\)) là số bộ
dữ liệu.

Dòng đầu tiên của mỗi bộ dữ liệu chứa hai số nguyên \(n\)
(\(1\le n\le 10^5\)) và \(k\)
(\(1\le k\le \frac{n(n+1)}{2}\)).

Dòng thứ hai của mỗi bộ dữ liệu chứa một chuỗi \(S\) độ dài \(n\), gồm các
chữ cái tiếng Anh thường.

Dòng thứ ba của mỗi bộ dữ liệu chứa \(26\) số nguyên
\(c_a,c_b,\ldots,c_z\) (\(1\le c_\alpha\le 100\)), là mức tiêu thụ năng lượng
của từng chữ cái.

Bảo đảm rằng tổng \(n\) trên tất cả các bộ dữ liệu không vượt quá
\(8\times 10^5\).

\subsection*{Dữ liệu ra}

Với mỗi bộ dữ liệu, in một số nguyên trên một dòng riêng: mức tiêu thụ năng
lượng nhỏ thứ \(k\), hoặc \(-1\) nếu không có đáp án.

\subsection*{Ví dụ}

\paragraph{standard input}
\begin{verbatim}
2
5 5
ababc
3 1 2 1 1 1 1 1 1 1 1 1 1 1 1 1 1 1 1 1 1 1 1 1 1 1
5 15
ababc
3 1 2 1 1 1 1 1 1 1 1 1 1 1 1 1 1 1 1 1 1 1 1 1 1 1
\end{verbatim}

\paragraph{standard output}
\begin{verbatim}
4
-1
\end{verbatim}

\section{Problem E. Didn't I Say to Make My Abilities Average in the Next Life?!}

\paragraph{Tệp vào:} standard input
\paragraph{Tệp ra:} standard output
\paragraph{Giới hạn bộ nhớ:} 256 megabytes

Khi chuyển sinh vào thế giới giả tưởng, Kurihara đã xin Đấng Sáng tạo ban cho
mình năng lực ở mức trung bình. Nhưng Đấng Sáng tạo lại rất kém toán, nên
ông ta xem ``trung bình'' là một nửa tổng của giá trị lớn nhất và nhỏ nhất
trong các giá trị.

Có \(n\) loại sinh vật trong thế giới giả tưởng, được đánh số từ \(1\) đến
\(n\). Mỗi sinh vật có một giá trị năng lực. Giá trị năng lực của loại sinh
vật thứ \(i\) là \(a_i\). Đấng Sáng tạo có \(m\) phương án ban năng lực. Với
phương án thứ \(i\), Đấng Sáng tạo sẽ chọn đều ngẫu nhiên một đoạn \([l,r]\)
(\(x\le l\le r\le y\)) từ một đoạn nào đó \([x,y]\)
(\(1\le x\le y\le n\)), tính ``trung bình'' của các giá trị năng lực từ sinh
vật thứ \(l\) đến sinh vật thứ \(r\) theo định nghĩa của ông ta, rồi ban giá
trị đó cho Kurihara. Chú ý rằng định nghĩa ``trung bình'' ở đây là một nửa
tổng của giá trị lớn nhất và nhỏ nhất trong các giá trị.

Kurihara muốn biết kỳ vọng toán học của giá trị năng lực mà cô sẽ được ban.

\subsection*{Dữ liệu vào}

Dòng đầu tiên chứa một số nguyên \(T\) (\(1\le T\le 10\)), biểu thị số bộ dữ
liệu.

Với mỗi bộ dữ liệu, dòng đầu tiên chứa hai số nguyên \(n,m\)
(\(1\le n,m\le 2\times 10^5\)), lần lượt biểu thị số loại sinh vật và số
phương án ban năng lực.

Dòng thứ hai chứa \(n\) số nguyên \(a_1,a_2,\ldots,a_n\)
(\(1\le a_i\le 10^9\)), biểu thị giá trị năng lực của từng sinh vật.

Trong \(m\) dòng tiếp theo, dòng thứ \(i\) (\(1\le i\le m\)) chứa hai số
nguyên \(x,y\) (\(1\le x\le y\le n\)), biểu thị phương án thứ \(i\).

Bảo đảm rằng với tất cả các bộ dữ liệu,
\[
  \sum n\le 3\times 10^5,\qquad
  \sum m\le 3\times 10^5.
\]

\subsection*{Dữ liệu ra}

Với mỗi bộ dữ liệu, in \(m\) dòng. Ở dòng thứ \(i\), in đáp án của phương án
thứ \(i\) dưới dạng phân số modulo \(10^9+7\) trên một dòng. Cụ thể, nếu đáp
án là \(\frac{P}{Q}\), bạn cần in \(P\cdot Q^{-1}\bmod (10^9+7)\), trong đó
\(Q^{-1}\) là nghịch đảo nhân của \(Q\) theo modulo \(10^9+7\).

\subsection*{Ví dụ}

\paragraph{standard input}
\begin{verbatim}
1
6 4
1 1 4 5 1 4
1 1
4 5
1 4
1 6
\end{verbatim}

\paragraph{standard output}
\begin{verbatim}
1
3
750000008
809523818
\end{verbatim}

\section{Problem F. Directed Minimum Spanning Tree}

\paragraph{Tệp vào:} standard input
\paragraph{Tệp ra:} standard output
\paragraph{Giới hạn bộ nhớ:} 256 megabytes

Yukikaze đang học lý thuyết đồ thị. Cô bị cuốn hút bởi một bài toán tối ưu
tổ hợp thú vị, gọi là bài toán cây khung nhỏ nhất có hướng.

Một đồ thị con \(T\) của đồ thị có hướng \(G=(V,E)\) được gọi là một cây
khung nhỏ nhất có hướng (DMST) gốc \(r\) nếu với mọi đỉnh \(v\in V-\{r\}\),
tồn tại đúng một đường đi trong \(T\) từ \(r\) đến \(v\). Trọng số của một
DMST là tổng trọng số các cạnh của nó.

Với mỗi đỉnh \(u\) trong đồ thị đã cho, Yukikaze muốn bạn tìm trọng số của
cây khung nhỏ nhất có hướng gốc tại đỉnh \(u\).

\subsection*{Dữ liệu vào}

Dữ liệu vào gồm nhiều bộ dữ liệu. Dòng đầu tiên chứa một số nguyên \(T\)
(\(1\le T\le 3000\)), biểu thị số bộ dữ liệu.

Với mỗi bộ dữ liệu, dòng đầu tiên chứa hai số nguyên \(n\)
(\(1\le n\le 10^5\)) và \(m\) (\(1\le m\le 2\times 10^5\)), lần lượt biểu thị
số đỉnh và số cạnh của đồ thị.

Mỗi dòng trong \(m\) dòng tiếp theo chứa ba số nguyên \(u_i,v_i,w_i\)
(\(1\le u_i,v_i\le n,\ 1\le w_i\le 10^9\)), biểu thị đỉnh đầu, đỉnh cuối và
trọng số của cạnh thứ \(i\). Chú ý rằng các cạnh có hướng.

Gọi \(S_n\) và \(S_m\) lần lượt là tổng các giá trị \(n\) và tổng các giá trị
\(m\) trong dữ liệu vào. Bảo đảm rằng
\[
  1\le S_n\le 5\times 10^5,\qquad
  1\le S_m\le 10^6.
\]

\subsection*{Dữ liệu ra}

Với mỗi bộ dữ liệu, in \(n\) dòng. Dòng thứ \(i\) chứa trọng số của cây khung
nhỏ nhất có hướng gốc tại đỉnh \(i\). Nếu không tồn tại DMST như vậy, in
\(-1\).

\subsection*{Ví dụ}

\paragraph{standard input}
\begin{verbatim}
2
5 6
1 2 5
2 3 6
3 1 7
1 4 3
4 5 4
5 1 2
6 8
1 4 1
2 5 7
4 3 4
5 6 12
6 2 9
3 5 10
3 1 23
4 2 17
\end{verbatim}

\paragraph{standard output}
\begin{verbatim}
18
20
19
17
16
36
-1
55
58
-1
-1
\end{verbatim}

\section{Problem G. Increasing Subsequence}

\paragraph{Tệp vào:} standard input
\paragraph{Tệp ra:} standard output
\paragraph{Giới hạn bộ nhớ:} 256 megabytes

Trong một dãy số nguyên \(a_1,a_2,\ldots,a_n\), nếu một dãy con tăng không
phải là dãy con của bất kỳ dãy con tăng nào khác, ta gọi nó là cực đại. Dãy
con là một dãy có thể thu được bằng cách xóa một số phần tử, có thể là không
xóa phần tử nào, khỏi dãy ban đầu.

Tìm hoặc đếm dãy con tăng dài nhất là một bài toán kinh điển. Bây giờ
Yukikaze muốn bạn đếm số dãy con tăng cực đại trong một số hoán vị theo
modulo \(998244353\). Một hoán vị độ dài \(n\) là một dãy số trong đó mỗi số
từ \(1\) đến \(n\) xuất hiện đúng một lần.

\subsection*{Dữ liệu vào}

Dòng đầu tiên chứa một số nguyên \(T\) (\(1\le T\le 10^4\)), biểu thị số bộ
dữ liệu.

Dòng đầu tiên của mỗi bộ dữ liệu chứa một số nguyên \(n\)
(\(1\le n\le 10^5\)), biểu thị độ dài của hoán vị.

Dòng thứ hai của mỗi bộ dữ liệu chứa \(n\) số nguyên \(a_1,a_2,\ldots,a_n\)
(\(1\le a_i\le n\)), biểu thị hoán vị. Bảo đảm rằng mỗi số từ \(1\) đến
\(n\) xuất hiện đúng một lần.

Tổng \(n\) trong tất cả các bộ dữ liệu không vượt quá \(2\times 10^5\).

\subsection*{Dữ liệu ra}

Với mỗi bộ dữ liệu, in một số nguyên biểu thị số dãy con tăng cực đại trong
hoán vị đã cho theo modulo \(998244353\).

\subsection*{Ví dụ}

\paragraph{standard input}
\begin{verbatim}
2
4
2 1 4 3
5
1 5 2 4 3
\end{verbatim}

\paragraph{standard output}
\begin{verbatim}
4
3
\end{verbatim}

\section{Problem H. Lawn of the Dead}

\paragraph{Tệp vào:} standard input
\paragraph{Tệp ra:} standard output
\paragraph{Giới hạn bộ nhớ:} 256 megabytes

Một ngày nọ, một thây ma đến Lawn of the Dead, có thể xem là một lưới
\(n\times m\). Ban đầu, anh ta đứng ở ô trên cùng bên trái, tức là \((1,1)\).

Vì bộ não của thây ma bị hỏng, anh ta không có khả năng định hướng tốt. Trong
một bước, anh ta chỉ có thể đi xuống từ \((i,j)\) đến \((i+1,j)\) hoặc đi
sang phải từ \((i,j)\) đến \((i,j+1)\).

Có \(k\) ``mìn lotato'' trên lưới. Quả mìn thứ \(i\) ở \((x_i,y_i)\). Để
tránh bị tiêu diệt, anh ta sẽ không bao giờ bước vào ô chứa ``mìn lotato''.

Vậy anh ta có thể đến được bao nhiêu ô? Tính cả ô xuất phát.

\subsection*{Dữ liệu vào}

Dòng đầu tiên chứa một số nguyên \(t\) (\(1\le t\le 20\)), biểu thị số bộ dữ
liệu.

Dòng đầu tiên của mỗi bộ dữ liệu chứa ba số nguyên \(n,m,k\)
(\(2\le n,m,k\le 10^5\)): có một lưới \(n\times m\), và có \(k\) ``mìn
lotato'' trên lưới.

Mỗi dòng trong \(k\) dòng tiếp theo chứa hai số nguyên \(x_i,y_i\)
(\(1\le x_i\le n,\ 1\le y_i\le m\)): có một ``mìn lotato'' tại \((x_i,y_i)\).
Bảo đảm rằng không có ``mìn lotato'' tại \((1,1)\) và không có hai quả mìn
nào ở cùng một ô.

Bảo đảm rằng
\[
  \sum n\le 7\cdot 10^5,\qquad
  \sum m\le 7\cdot 10^5.
\]

\subsection*{Dữ liệu ra}

Với mỗi bộ dữ liệu, in số ô mà anh ta có thể đến được.

\subsection*{Ví dụ}

\paragraph{standard input}
\begin{verbatim}
1
4 4 4
1 3
3 4
3 2
4 3
\end{verbatim}

\paragraph{standard output}
\begin{verbatim}
10
\end{verbatim}

\subsection*{Ghi chú}

Các ô mà thây ma có thể đến là \((1,1)\), \((1,2)\), \((2,1)\), \((2,2)\),
\((2,3)\), \((2,4)\), \((3,1)\), \((3,3)\), \((4,1)\), \((4,2)\).

\section{Problem I. License Plate Recognition}

\paragraph{Tệp vào:} standard input
\paragraph{Tệp ra:} standard output
\paragraph{Giới hạn bộ nhớ:} 256 megabytes

Little Rabbit là một sinh viên đại học đang học Kỹ thuật Truyền thông. Trong
học kỳ này, cậu có một môn học tên là Xử lý ảnh số. Cuối môn học, Little
Rabbit cần hoàn thành một dự án tên là Nhận dạng biển số xe, trong đó Little
Rabbit phải lập trình một hệ thống nhận dạng các ký tự của biển số xe trong
một ảnh.

Một hệ thống nhận dạng biển số xe kinh điển có ba mô-đun:
\begin{itemize}
  \item Định vị biển số xe: định vị biển số xe trong ảnh.
  \item Phân đoạn ký tự: phân đoạn ảnh để các ký tự của biển số xe có thể
  được tách ra.
  \item Nhận dạng ký tự: nhận dạng các ký tự trong ảnh và nhận chuỗi kết quả.
\end{itemize}

Để hoàn thành dự án, Little Rabbit lập một đội gồm ba thành viên. Mỗi thành
viên phụ trách một mô-đun. Ở đây, Little Rabbit phụ trách mô-đun thứ hai:
Phân đoạn ký tự.

Sau mô-đun Định vị biển số xe và một số bước tiền xử lý, Little Rabbit nhận
được một số ảnh nhị phân biểu diễn biển số xe. Tất cả các ảnh nhị phân đều có
chiều rộng \(100\) pixel và chiều cao \(30\) pixel, chỉ chứa pixel đen và
pixel trắng. Pixel đen biểu diễn nền, còn pixel trắng biểu diễn các ký tự.

Nhiệm vụ của Little Rabbit là phân đoạn các ảnh nhị phân để các ký tự trong
ảnh có thể được tách ra. Theo tiêu chuẩn, có bảy ký tự trong một biển số xe,
xếp từ trái sang phải. Cụ thể, nhiệm vụ của Little Rabbit là tìm biên trái và
biên phải của mỗi ký tự.

Hãy xem ảnh như các ma trận có \(30\) hàng và \(100\) cột. Đánh số các cột từ
\(1\) đến \(100\) theo chiều từ trái sang phải. Ta định nghĩa biên trái của
một ký tự là chỉ số cột chứa pixel ngoài cùng bên trái của ký tự đó, và biên
phải là chỉ số cột chứa pixel ngoài cùng bên phải. Ví dụ, trong hình ở đề
gốc, biên trái của ký tự là \(3\), và biên phải là \(7\).

Cho các ảnh nhị phân mà Little Rabbit cần phân đoạn, hãy in biên trái và biên
phải của từng ký tự. Trong bài này, ta dùng \texttt{.} để biểu diễn một pixel
đen, và \texttt{\#} để biểu diễn một pixel trắng.

\subsection*{Dữ liệu vào}

Dòng đầu tiên chứa một số nguyên \(T\) (\(1\le T\le 50\)) là số bộ dữ liệu.

Mỗi bộ dữ liệu biểu diễn một ảnh nhị phân, gồm \(30\) dòng chuỗi. Mỗi dòng
chứa \(100\) ký tự, là \texttt{.} hoặc \texttt{\#}.

Sau đây là tất cả các ký tự có thể xuất hiện trong ảnh.

Ký tự tiếng Trung:

Phiên bản ASCII: \url{https://paste.ubuntu.com/p/B5pTWv7s6J/}

Bản dự phòng: \url{https://github.com/cjj490168650/plate/blob/main/chn.txt}

Ký tự tiếng Anh và chữ số:

Phiên bản ASCII: \url{https://paste.ubuntu.com/p/59bjvwY3Yr/}

Bản dự phòng: \url{https://github.com/cjj490168650/plate/blob/main/eng.txt}

Bảo đảm rằng:
\begin{enumerate}
  \item Các ký tự trong ảnh tuân theo tiêu chuẩn của biển số xe. Có bảy ký tự
  trong ảnh, xếp từ trái sang phải. Ký tự đầu tiên là một ký tự tiếng Trung.
  Ký tự thứ hai là một ký tự tiếng Anh. Năm ký tự cuối là ký tự tiếng Anh
  hoặc chữ số.
  \item Tất cả các ký tự trong ảnh giống hệt các ký tự đã cho ở trên, gồm cả
  kích thước và hình dạng trong phiên bản ASCII.
  \item Không có pixel trắng thừa trong ảnh.
  \item Có khoảng cách giữa các ký tự.
  \item Các ký tự sẽ không chạm vào nhau và không vượt ra ngoài biên ảnh.
\end{enumerate}

\subsection*{Dữ liệu ra}

Với bộ dữ liệu thứ \(x\), in \texttt{Case \#x:} ở dòng đầu tiên.

Sau đó, trong bảy dòng tiếp theo, dòng thứ \(i\) chứa hai số nguyên cách nhau
bởi một dấu cách, biểu thị biên trái và biên phải của ký tự thứ \(i\).

\subsection*{Ghi chú}

Dữ liệu vào và dữ liệu ra mẫu:
\url{https://paste.ubuntu.com/p/28rYFsbKHb/}

Bản dự phòng:
\url{https://github.com/cjj490168650/plate/blob/main/sample.txt}

\section{Problem J. Pony Running}

\paragraph{Tệp vào:} standard input
\paragraph{Tệp ra:} standard output
\paragraph{Giới hạn bộ nhớ:} 256 megabytes

Ponyville đã bị changeling kiểm soát! Các pony quá sợ hãi nên không dám di
chuyển. May mắn là trong Ponyville vẫn còn một vùng phép thuật để bảo vệ các
pony.

Vùng phép thuật là một lưới \(n\times m\). Ta ký hiệu ô ở hàng thứ \(x\) và
cột thứ \(y\) là \((x,y)\). Mỗi giây, các pony có thể di chuyển lên, xuống,
trái hoặc phải với các xác suất khác nhau. Cụ thể, nếu một pony đang ở
\((x,y)\), thì trong giây tiếp theo, nó sẽ đi lên \((x-1,y)\) với xác suất
\(p_{xy0}\), đi xuống \((x+1,y)\) với xác suất \(p_{xy1}\), đi sang trái
\((x,y-1)\) với xác suất \(p_{xy2}\), hoặc đi sang phải \((x,y+1)\) với xác
suất \(p_{xy3}\). Các pony có thể đi ra ngoài vùng phép thuật, chẳng hạn một
pony ở ô thuộc hàng đầu tiên và đi lên trong giây tiếp theo. Vì việc đi ra
ngoài vùng phép thuật rất nguy hiểm, các pony muốn biết tổng kỳ vọng toán học
của thời gian để mỗi pony đi ra khỏi vùng.

Có \(q\) sự kiện. Loại sự kiện thứ nhất là thay đổi xác suất cho bốn hướng
tại một ô. Loại sự kiện thứ hai là đặt một pony ở mỗi ô, rồi truy vấn tổng kỳ
vọng toán học của thời gian để mỗi pony đi ra khỏi vùng.

\subsection*{Dữ liệu vào}

Dòng đầu tiên chứa ba số nguyên \(n,m,q\)
(\(1\le n,m,q\le 400,\ 1\le n\times m\le 400\)), biểu thị số hàng, số cột của
lưới và số sự kiện.

Trong \(n\times m\) dòng tiếp theo, dòng thứ \(((i-1)\times m+j)\) chứa
\(4\) số nguyên \(p_{ij0},p_{ij1},p_{ij2},p_{ij3}\)
(\(0\le p_{ij0},p_{ij1},p_{ij2},p_{ij3}<1{,}000{,}000{,}007\)), biểu thị xác
suất để pony ở \((i,j)\) đi lên, xuống, trái và phải. Các xác suất được cho
dưới dạng modulo \(1{,}000{,}000{,}007\). Cụ thể, nếu xác suất là
\(P/Q\), thì số nguyên được cho là
\(P\cdot Q^{-1}\pmod {1{,}000{,}000{,}007}\), trong đó \(Q^{-1}\) là nghịch
đảo nhân của \(Q\) theo modulo \(1{,}000{,}000{,}007\). Bảo đảm rằng
\[
  p_{ij0}+p_{ij1}+p_{ij2}+p_{ij3}\equiv 1
  \pmod {1{,}000{,}000{,}007}.
\]

Trong \(q\) dòng tiếp theo, mỗi dòng biểu diễn một sự kiện.
\begin{itemize}
  \item \texttt{1 x y p0 p1 p2 p3}, biểu thị bốn xác suất của ô \((x,y)\) sẽ
  được đổi thành \(p_0,p_1,p_2,p_3\).
  \[
    1\le x\le n,\quad 1\le y\le m,\quad
    0\le p_0,p_1,p_2,p_3<1{,}000{,}000{,}007,
  \]
  \[
    p_0+p_1+p_2+p_3\equiv 1\pmod {1{,}000{,}000{,}007}.
  \]
  \item \texttt{2}, biểu thị truy vấn tổng kỳ vọng toán học của thời gian để
  mỗi pony đi ra khỏi vùng.
\end{itemize}

\subsection*{Dữ liệu ra}

Với mỗi truy vấn, in đáp án theo modulo \(1{,}000{,}000{,}007\) trên một
dòng. Cụ thể, nếu đáp án là \(P/Q\), bạn cần in
\(P\cdot Q^{-1}\pmod {1{,}000{,}000{,}007}\), trong đó \(Q^{-1}\) là nghịch
đảo nhân của \(Q\) theo modulo \(1{,}000{,}000{,}007\).

\subsection*{Ví dụ}

\paragraph{standard input}
\begin{verbatim}
2 2 5
500000004 0 500000004 0
500000004 0 500000004 0
500000004 0 500000004 0
500000004 0 500000004 0
2
1 1 1 0 500000004 0 500000004
2
1 1 2 0 500000004 0 500000004
2
\end{verbatim}

\paragraph{standard output}
\begin{verbatim}
500000010
14
14
\end{verbatim}

\section{Problem K. Travel on Tree}

\paragraph{Tệp vào:} standard input
\paragraph{Tệp ra:} standard output
\paragraph{Giới hạn bộ nhớ:} 256 megabytes

Trong thế giới của The Three-Body Problem, khoảng \(200\) năm sau, con người
sẽ sống trong các công trình cây khổng lồ dưới lòng đất. Trong bài này, ta
dùng cấu trúc dữ liệu cây để mô tả các công trình cây.

Có một cây gồm \(n\) nút và \(n-1\) cạnh có độ dài, và \(n\) người bạn của
bạn sống trong cây, mỗi người ở một nút. Bạn dự định thăm bạn bè trong
\(m\) ngày tiếp theo. Mỗi ngày bạn chọn một đoạn \([l,r]\) và dự định thăm
những người bạn sống ở các nút được đánh số từ \(l\) đến \(r\). Bạn có thể
chọn một nút bất kỳ \(u\) để bắt đầu chuyến thăm trong ngày, sau đó di chuyển
trên cây dọc theo các cạnh, và cuối cùng quay lại \(u\). Trong chuyến đi, bạn
phải thăm tất cả bạn bè sống ở các nút được đánh số từ \(l\) đến \(r\). Bạn
có thể thăm những người bạn này theo bất kỳ thứ tự nào và có thể đi qua một
nút mà không thăm người bạn ở đó. Hãy tính tổng khoảng cách di chuyển nhỏ
nhất cho mỗi ngày.

\subsection*{Dữ liệu vào}

Dòng đầu tiên chứa một số nguyên \(T\) (\(1\le T\le 2\times 10^4\)) là số bộ
dữ liệu.

Dòng đầu tiên của mỗi bộ dữ liệu chứa hai số nguyên \(n,m\)
(\(1\le n,m\le 10^5\)): số nút và số ngày.

Mỗi dòng trong \(n-1\) dòng tiếp theo của mỗi bộ dữ liệu chứa ba số nguyên
\(u,v,w\) (\(1\le u,v\le n,\ 1\le w\le 10^4\)): một cạnh nối các nút \(u\) và
\(v\) có độ dài \(w\).

Mỗi dòng trong \(m\) dòng tiếp theo của mỗi bộ dữ liệu chứa hai số nguyên
\(l,r\) (\(1\le l\le r\le n\)): một kế hoạch thăm những người bạn sống ở các
nút được đánh số từ \(l\) đến \(r\).

Bảo đảm rằng với tất cả các bộ dữ liệu,
\[
  \sum n\le 10^6,\qquad
  \sum m\le 10^6.
\]

\subsection*{Dữ liệu ra}

Với kế hoạch của mỗi ngày, in đáp án trên một dòng.

\subsection*{Ví dụ}

\paragraph{standard input}
\begin{verbatim}
2
3 3
1 2 1
2 3 1
1 1
1 2
1 3
5 5
1 2 1
1 3 2
3 4 3
3 5 4
1 3
1 4
1 5
2 5
3 5
\end{verbatim}

\paragraph{standard output}
\begin{verbatim}
0
2
4
6
12
20
20
14
\end{verbatim}

\end{document}
